\chapter{Технологическая часть}

В данном разделе описаны инструменты, использованные для реализации, а также представлены листинги трёх алгоритмов умножения матриц: стандартного, классического алгоритма Винограда и его оптимизированной версии.

\section{Средства реализации}
Для реализации был выбран язык C++ в качестве среды разработки использовалась Visual Studio Code с компилятором g++. Для модульного тестирования применялась библиотека Check, а для визуализации результатов~--- утилита gnuplot.

\section{Реализация алгоритмов}

В~листингах~\ref{lst:standard}~---\ref{lst:opt_vinograd} представлены реализации алгоритмов умножения матриц.

\section{Функциональные тесты}

Для проверки корректности реализации алгоритмов умножения матриц были разработаны функциональные тесты, охватывающие различные классы эквивалентности. Результаты тестирования представлены в таблице~\ref{tbl:functional_test}.

\begin{table}[h]
	\centering
	\begin{threeparttable}
		\caption{Функциональные тесты алгоритмов умножения матриц}
		\label{tbl:functional_test}
		\begin{tabular}{|p{0.25\textwidth}|p{0.25\textwidth}|p{0.3\textwidth}|}
			\hline
			\textbf{Входная матрица A} & \textbf{Входная матрица B} & \textbf{Ожидаемый результат} \\ 
			\hline
			
			$\begin{pmatrix}
				1 & 2 \\
				3 & 4
			\end{pmatrix}$ &
			$\begin{pmatrix}
				5 & 6 \\
				7 & 8
			\end{pmatrix}$ &
			$\begin{pmatrix}
				19 & 22 \\
				43 & 50
			\end{pmatrix}$ \\
			\hline
			
			$\begin{pmatrix}
				1 & 2 \\
				3 & 4 \\
				5 & 6
			\end{pmatrix}$ &
			$\begin{pmatrix}
				7 & 8 & 9 \\
				10 & 11 & 12
			\end{pmatrix}$ &
			$\begin{pmatrix}
				27 & 30 & 33 \\
				61 & 68 & 75 \\
				95 & 106 & 117
			\end{pmatrix}$ \\
			\hline
			
			$\begin{pmatrix}
				1 & 2 \\
				3 & 4
			\end{pmatrix}$ &
			$\begin{pmatrix}
				1 & 0 \\
				0 & 1
			\end{pmatrix}$ &
			$\begin{pmatrix}
				1 & 2 \\
				3 & 4
			\end{pmatrix}$ \\
			\hline
			
			$\begin{pmatrix}
				1 & 2 \\
				3 & 4
			\end{pmatrix}$ &
			$\begin{pmatrix}
				0 & 0 \\
				0 & 0
			\end{pmatrix}$ &
			$\begin{pmatrix}
				0 & 0 \\
				0 & 0
			\end{pmatrix}$ \\
			\hline
			
			$\begin{pmatrix}
				5
			\end{pmatrix}$ &
			$\begin{pmatrix}
				7
			\end{pmatrix}$ &
			$\begin{pmatrix}
				35
			\end{pmatrix}$ \\
			\hline
			
			$\begin{pmatrix}
				1 & 2 & 3
			\end{pmatrix}$ &
			$\begin{pmatrix}
				4 \\
				5 \\
				6
			\end{pmatrix}$ &
			$\begin{pmatrix}
				32
			\end{pmatrix}$ \\
			\hline
			
			$\begin{pmatrix}
				1 \\
				2 \\
				3
			\end{pmatrix}$ &
			$\begin{pmatrix}
				4 & 5 & 6
			\end{pmatrix}$ &
			$\begin{pmatrix}
				4 & 5 & 6 \\
				8 & 10 & 12 \\
				12 & 15 & 18
			\end{pmatrix}$ \\
			\hline
			
			$\begin{pmatrix}
				1 & 2 & 3
			\end{pmatrix}$ &
			$\begin{pmatrix}
				1 & 2 \\
				3 & 4
			\end{pmatrix}$ &
			Сообщение об ошибке (несовместимые размеры) \\
			\hline
		\end{tabular}
	\end{threeparttable}
\end{table}

Разработанные тесты охватывают следующие классы эквивалентности:
\begin{enumerate}
	\item умножение квадратных матриц одинаковой размерности;
	\item умножение прямоугольных матриц различных размеров;
	\item умножение на единичную матрицу (проверка нейтрального элемента);
	\item умножение на нулевую матрицу (проверка поглощающего элемента);
	\item умножение матриц-векторов (граничные случаи размерности);
	\item обработка несовместимых размеров матриц;
\end{enumerate}

Все тестовые случаи успешно пройдены, реализация всех трёх алгоритмов корректна.

\section{Вывод}

В технологической части представлены листинги реализаций алгоритмов умножения матриц, описаны выбранные средства разработки и приведены функциональные тесты. Тестирование подтвердило корректность работы всех реализованных алгоритмов для различных классов входных данных.

\clearpage